% RQ3 Conclusion: Proposal Pipeline Effects

\subsection{Summary}

We have presented a systematic analysis of proposal pipeline design through multi-agent simulation. Across \experimentruns{} simulation runs testing configurations of temp-checks, fast-tracks, and expiry windows, we characterized the throughput-quality tradeoffs inherent in pipeline design.

Key findings:

\begin{enumerate}
    \item \textbf{Temp-checks provide valuable filtering}: Moderate thresholds (20-30\%) reduce formal voting load while maintaining proposal quality. Higher thresholds filter too aggressively.

    \item \textbf{Fast-tracks accelerate consensus}: 70\% support threshold effectively identifies uncontroversial proposals suitable for expedited processing. Higher thresholds are too restrictive.

    \item \textbf{Expiry windows should match complexity}: 30-day windows cause excessive abandonment; 60-day windows provide better balance. Complex proposals may need longer.

    \item \textbf{Moderate configurations dominate}: Both extremes (no filtering vs. aggressive filtering) underperform moderate, balanced approaches.
\end{enumerate}

\subsection{Contributions}

This paper contributes:

\begin{enumerate}
    \item \textbf{Formal pipeline models} with explicit throughput and quality metrics
    \item \textbf{Systematic parameter analysis} for three key pipeline mechanisms
    \item \textbf{Design principles} for pipeline configuration
    \item \textbf{Parameter guidelines} based on simulation evidence
\end{enumerate}

\subsection{Practical Recommendations}

For DAO practitioners designing proposal pipelines:

\begin{enumerate}
    \item \textbf{Implement temp-checks}: Off-chain signal votes reduce on-chain voting costs and filter low-quality proposals
    \item \textbf{Set moderate thresholds}: 20-30\% for temp-checks, 70\% for fast-tracks
    \item \textbf{Match expiry to context}: 60 days for standard proposals; consider longer for complex governance
    \item \textbf{Monitor metrics}: Track abandonment, expiry, and false negative rates to tune parameters
\end{enumerate}

\subsection{Implications}

For governance designers, our results demonstrate that pipeline design significantly impacts governance outcomes. The choice is not whether to have a pipeline but how to configure it for the DAO's specific needs.

For researchers, our framework enables continued investigation of pipeline dynamics. The interaction between pipeline parameters and other governance mechanisms (voting rules, delegation) deserves further study.

\subsection{Closing Remarks}

Proposal pipelines are the operational backbone of DAO governance. A well-designed pipeline balances efficiency (processing proposals quickly) with quality (ensuring adequate deliberation). Our simulations reveal that this balance is achievable through moderate, evidence-based parameter selection.

We release our simulation framework to enable the community to extend this analysis and develop best practices for pipeline design across diverse DAO contexts.
